\section{Scaling laws and sensitivity}
\label{sec:scaling}

\Cref{cor:closedform} gives the critical payload in closed form. Because it is
a product of powers of the underlying parameters, its logarithmic derivatives
are exact rational numbers, and those numbers are the design guidance.

\subsection{The explicit critical payload}

\begin{proposition}[Explicit form of $\Mmax$]
\label{prop:mmaxexplicit}
For $q=3/2$,
\begin{equation}
  \Mmax
  = \frac{4\alpha^{3}}{27\beta^{2}}
  = \frac{8\,\alpha^{3}\,\rho\, \Adisk\, \eb^{2}\,\FM^{2}\,\eta^{2}}
         {27\, g^{3}\, t_f^{2}} .
  \label{eq:mmaxexplicit}
\end{equation}
\end{proposition}

\begin{proof}
From \cref{eq:beta}, $\beta^2 = t_f^2 g^3/(\eb^2\FM^2\eta^2\cdot 2\rho A)$.
Hence $1/\beta^2 = 2\rho A \eb^2\FM^2\eta^2/(t_f^2 g^3)$ and
$\Mmax = (4\alpha^3/27)\cdot 2\rho A\eb^2\FM^2\eta^2/(t_f^2g^3)$, which is
\cref{eq:mmaxexplicit}.
\end{proof}

\subsection{Exact logarithmic sensitivities}

\begin{definition}
For a positive quantity $Y$ depending on a positive parameter $x$, the
logarithmic sensitivity is
$S_x \coloneqq \partial \ln Y/\partial \ln x$, so that a relative change
$\delta x/x$ produces $S_x\,\delta x/x$ to first order.
\end{definition}

\begin{theorem}[Sensitivities of the critical payload]
\label{thm:sensitivity}
With $\Mmax$ as in \cref{eq:mmaxexplicit} and
$\gamma_m = \lambda g/S_m$, $\alpha = 1-f_s-\gamma_m$,
\begin{align}
  S_{t_f} &= -2, &
  S_{\eb} &= +2, &
  S_{\FM} &= +2, &
  S_{\eta} &= +2, \label{eq:sens1}\\
  S_{\Adisk} &= +1, &
  S_{D} &= +2, &
  S_{N} &= +1, &
  S_{\rho} &= +1, \label{eq:sens2}
\end{align}
and for the parameters entering through $\alpha$,
\begin{align}
  S_{f_s} &= -\frac{3 f_s}{\alpha}, &
  S_{\lambda} &= -\frac{3\gamma_m}{\alpha}, &
  S_{S_m} &= +\frac{3\gamma_m}{\alpha}, &
  S_{g} &= -3 - \frac{3\gamma_m}{\alpha} . \label{eq:sens3}
\end{align}
\end{theorem}

\begin{proof}
Take logarithms of \cref{eq:mmaxexplicit}:
\[
  \ln \Mmax = \ln\tfrac{8}{27} + 3\ln\alpha + \ln\rho + \ln A
              + 2\ln \eb + 2\ln\FM + 2\ln\eta - 3\ln g - 2\ln t_f .
\]
For any parameter $x$ not appearing in $\alpha$, differentiation is immediate
and gives \cref{eq:sens1,eq:sens2}; the entries for $D$ and $N$ follow from
$A = N\pi D^2/4$, so $\partial\ln A/\partial\ln D = 2$ and
$\partial\ln A/\partial \ln N = 1$.

For $f_s$: $\partial\alpha/\partial f_s = -1$, so
$\partial \ln\Mmax/\partial\ln f_s = 3 f_s\,\partial\ln\alpha/\partial f_s
= -3f_s/\alpha$.

For $\lambda$: $\alpha = 1-f_s-\lambda g/S_m$, so
$\partial\alpha/\partial\lambda = -g/S_m$ and
$\partial\ln\Mmax/\partial\ln\lambda
= 3\lambda(-g/S_m)/\alpha = -3\gamma_m/\alpha$. The result for $S_m$ follows
identically with the opposite sign since $\gamma_m\propto S_m^{-1}$.

For $g$: gravity appears both explicitly, as $g^{-3}$, and inside $\alpha$
through $\gamma_m = \lambda g/S_m$. Hence
\[
  S_g = -3 + 3\frac{g}{\alpha}\frac{\partial \alpha}{\partial g}
      = -3 + \frac{3g}{\alpha}\left(-\frac{\lambda}{S_m}\right)
      = -3 - \frac{3\gamma_m}{\alpha}. \qedhere
\]
\end{proof}

\begin{corollary}[The inverse-square endurance law]
\label{cor:tsquared}
At fixed technology and geometry, $\Mmax \propto t_f^{-2}$. Multiplying the
required endurance by $k$ divides the supportable fixed load by $k^2$.
\end{corollary}

\Cref{fig:walltf} shows the law for the lumped model of the design of
\cref{sec:results}, and what it implies once the auxiliary load is counted.

\begin{figure}[tbp]
\centering
% Lumped fixed-area closure of the design of sec17 (gen_b.wall_tf):
% the critical load M0max(t_f) falls as t_f^{-2}; the load the mission demands,
% M0(t_f) = m_fix + P_aux t_f / e_b, rises. They cross at the endurance wall.
\begin{tikzpicture}
\begin{axis}[paperplot, height=6.2cm, xmode=log, ymode=log,
  xlabel={mission duration $t_f$ [\si{\minute}]},
  ylabel={load [\si{\kilo\gram}]},
  xmin=5, xmax=480, ymin=0.01, ymax=150, legend pos=north east,
  xtick={5,10,20,30,60,120,240,480},
  xticklabels={5,10,20,30,60,120,240,480}]
  \addplot[cA, name path=wall] table[x=tmin, y=Mmax] {results/fig/b_walltf.dat};
  \addlegendentry{critical load $\Mmax(t_f)\propto t_f^{-2}$}
  \addplot[cB, name path=need] table[x=tmin, y=M0] {results/fig/b_walltf.dat};
  \addlegendentry{required load $\Mzero(t_f)=\mfix+\Paux t_f/\eb$}
  \path[name intersections={of=wall and need, name=w}];
  \fill[black] (w-1) circle (2.2pt);
  \draw[construct] (w-1) -- (w-1 |- {axis cs:5,0.01});
  \node[lbl, anchor=south west, align=left] at (w-1) {endurance wall\\of the lumped\\fixed-area closure};
  % slope triangle
  \draw[cG] (axis cs:10,0.5) -- (axis cs:40,0.5) -- (axis cs:10,8) -- cycle;
  \node[lbl, cG, anchor=north] at (axis cs:20,0.5) {$\times4$};
  \node[lbl, cG, anchor=east] at (axis cs:10,2) {$\times16$};
  \node[lbl, cG, anchor=west] at (axis cs:10.5,1.2) {slope $-2$};
  \draw[cG, densely dotted, line width=0.7pt] (axis cs:30,0.01) -- (axis cs:30,150);
  \node[lbl, cG, anchor=west] at (axis cs:30,0.02) {design $t_f$};
\end{axis}
\end{tikzpicture}

\caption{The inverse-square law of \cref{cor:tsquared} for the lumped
fixed-area model of the design of \cref{sec:results}. The critical load
$\Mmax$ falls with slope $-2$ on logarithmic axes. The load the mission itself
demands rises, because the battery that runs the auxiliary load,
$\Paux t_f/\eb$, is part of $\Mzero$. The two cross at the endurance wall of
this closure, \SI{66}{\minute}, against the design's \SI{30}{\minute} session;
beyond the crossing no mass closes.}
\label{fig:walltf}
\end{figure}

\begin{example}[Verification against the engine]
For the parameter set of \cref{sec:results} the engine computes the
sensitivities numerically by central differences in $\ln$-space and returns
$S_{t_f}=-2.000$, $S_{\eb}=S_{\FM}=S_{\eta}=+2.000$, $S_D=+2.000$, $S_N=+1.000$,
$S_\rho=+1.000$, $S_{f_s}=-0.943$, $S_{S_m}=+0.770$, $S_\lambda=-0.770$ and
$S_g=-3.770$. With $\gamma_m = 2\times 9.80665/120 = 0.16344$ and
$\alpha=0.63656$, \cref{eq:sens3} predicts
$S_{f_s} = -3(0.20)/0.63656 = -0.9426$,
$S_{S_m} = 3(0.16344)/0.63656 = 0.7703$ and
$S_g = -3-0.7703 = -3.7703$. The agreement to the printed precision is a
verification of both the closed form and its implementation. See
\cref{sec:verification}.
\end{example}

\begin{remark}[How to read \cref{thm:sensitivity}]
Three groups of parameters behave quite differently.

The \emph{squared} levers, $\eb$, $\FM$, $\eta$ and $t_f$, are the ones with
real authority, and their sign structure is unforgiving: a 20\% better cell
buys 44\% more payload, once; a doubling of endurance costs three quarters of
it. The figure of merit and the electrical efficiency are bounded above by
unity, so at $\FM=0.73$ and $\eta=0.78$ there remains at most a factor
$(1/0.73)^2(1/0.78)^2 \approx 3.1$ in that direction, ever.

The \emph{linear} levers are disk area and air density. Disk area is the only
one of the four that the designer controls freely, and because $D$ enters
squared it is in practice the strongest available lever.

The \emph{cubic-in-$\alpha$} levers, $f_s$, $\lambda$, $S_m$, have modest
sensitivity as long as $\alpha$ is comfortably positive, and become violent as
$\alpha\to0$. They are not where the design is won or lost unless one has been
careless.
\end{remark}

\Cref{fig:sensitivity} collects the sensitivities as the engine computes them.

\begin{figure}[tbp]
\centering
% Logarithmic sensitivities of M0max computed by the engine
% (sizing.sensitivities at the default parameters, central differences in
% log space). Order written by gen_b.sensitivities.
\begin{tikzpicture}
\begin{axis}[paperplot, height=6.6cm, width=0.8\linewidth,
  xbar, bar width=7.5pt, bar shift=0pt, y dir=reverse,
  ytick={0,...,10},
  yticklabels={$t_f$,$\eb$,$\FM$,$\eta$,$D$,$N$,$\rho$,$f_s$,$\lambda$,$S_m$,$g$},
  xmin=-4.6, xmax=3.1, ymin=-0.6, ymax=10.6,
  xlabel={$S_x=\partial\ln\Mmax/\partial\ln x$},
  grid=none, xmajorgrids, xtick={-4,-3,...,3},
  nodes near coords, point meta=x,
  every node near coord/.append style={font=\scriptsize,
    /pgf/number format/fixed, /pgf/number format/precision=2,
    /pgf/number format/fixed zerofill}]
  \addplot[fill=cA!65, draw=cA] table[x=S, y=i, restrict x to domain=0:10]
    {results/fig/b_sens.dat};
  \addplot[fill=cD!55, draw=cD] table[x=S, y=i, restrict x to domain=-10:0]
    {results/fig/b_sens.dat};
  \draw[cG] (axis cs:0,-0.6) -- (axis cs:0,10.6);
  \draw[cG, dashed] (axis cs:-4.6,3.5) -- (axis cs:3.1,3.5);
  \draw[cG, dashed] (axis cs:-4.6,6.5) -- (axis cs:3.1,6.5);
  \node[lbl, cG, anchor=west] at (axis cs:-4.5,1.5) {squared};
  \node[lbl, cG, anchor=west] at (axis cs:-4.5,5.0) {disk area, air};
  \node[lbl, cG, anchor=west, align=left] at (axis cs:-4.5,7.6) {through $\alpha$};
\end{axis}
\end{tikzpicture}

\caption{Logarithmic sensitivities of $\Mmax$ at the engine's default
parameters, computed by central differences in $\ln$-space. The top group are
the squared levers of \cref{eq:sens1}; the middle group are disk area, entering
as $D^2$ and $N$, and air density, \cref{eq:sens2}; the bottom group act
through $\alpha$ and are \cref{eq:sens3}, here with $\gamma_m/\alpha=0.257$.
Every bar agrees with \cref{thm:sensitivity} to three decimals. A bar of length
two means a \SI{10}{\percent} change in the parameter moves the critical load
by about \SI{20}{\percent}.}
\label{fig:sensitivity}
\end{figure}

\subsection{The exponent, the constants, and where the wall sits}

A claim one meets in discussions of this problem is that ``the exponents kill
you, not the constants''. \Cref{thm:sensitivity} lets us be precise about the
sense in which that is and is not true.

\begin{observation}
\label{obs:expconst}
The exponent $q$ determines whether a wall exists at all
(\cref{thm:pfold}: it exists iff $q>1$) and the asymptotic behaviour of the
closure. The constants $\alpha$ and $\beta$ determine where the wall sits,
through \cref{eq:closedform}, and their influence is not weak: $\Mmax$ is
cubic in $\alpha$ and inverse-quadratic in $\beta$. The correct statement is
that the exponent creates the wall and the constants locate it.
\end{observation}

For finite engineering at finite scales the location is what one actually
cares about. \Cref{eq:mmaxexplicit} says that location moves by a factor of two
for a 41\% improvement in cell energy density, or a 41\% improvement in figure
of merit, or a doubling of disk area, or a 29\% reduction in endurance.

\subsection{The endurance wall across the exponent}

Combining \cref{thm:pfold} with \cref{prop:plaw} gives the full picture, but
only once two statements that are easily conflated are separated: whether the
closure has a fold, and whether there is a longest mission it can close. They
are different, and they part company at $p=1$.

Write the closure with its dependence on duration explicit. Under
\cref{eq:areascaling} the propulsive power is $c_P m^q$ with
$c_P = c_0 m_0^{p/2}$ and $c_0 = g^{3/2}/(\FM\,\eta\sqrt{2\rho A_0})$, so a mass
$m$ closes a mission of duration $t$ if and only if
$\alpha m - (t/\eb)(c_P m^q + \Paux) = \mfix$, that is if and only if
$t = \Theta_p(m)$ with
\begin{equation}
  \Theta_p(m) \coloneqq
  \frac{\eb\,(\alpha m - \mfix)}{c_P\, m^{q} + \Paux} .
  \label{eq:theta}
\end{equation}
Define the endurance wall at fixed $\mfix$, $\Paux$ and anchor $(m_0,A_0)$ as
$t_f^{\max}(p) \coloneqq \sup_{m>0}\Theta_p(m)$.

\begin{proposition}[The endurance wall]
\label{prop:twall}
Let $\alpha>0$ and $\mfix>0$.
\begin{enumerate}[label=(\roman*),nosep]
  \item If $p<1$, $\Theta_p$ attains its maximum at a finite mass
        $\hat m(p) > \mfix/\alpha$, and $t_f^{\max}(p)=\Theta_p(\hat m)$ is
        finite and attained.
  \item If $p=1$, $\Theta_1$ is strictly increasing and
        \begin{equation}
          t_f^{\max}(1) = t_{\lim} = \frac{\alpha\,\eb}{c_P} ,
          \label{eq:tlim}
        \end{equation}
        which is finite and not attained: for every $t_f<t_{\lim}$ the unique
        mass is finite, and it diverges as $t_f\uparrow t_{\lim}$. Here $\eb$
        is net of the landing reserve, as in \cref{mod:battery}.
  \item If $p>1$, $\Theta_p(m)\to\infty$ as $m\to\infty$ and there is no wall.
  \item Where $\hat m$ is differentiable in $p$,
        \begin{equation}
          \frac{\dd\ln t_f^{\max}}{\dd p}
          = \frac{\vartheta}{2}\,\ln\frac{\hat m}{m_0} ,
          \qquad
          \vartheta = \frac{c_P\hat m^{q}}{c_P \hat m^{q}+\Paux}\in(0,1],
          \label{eq:wallslope}
        \end{equation}
        so the wall rises with $p$ if and only if the mass at the wall exceeds
        the anchor mass. It is not monotone in $p$ for every anchor.
\end{enumerate}
\end{proposition}

\begin{proof}
$\Theta_p$ is continuous, negative for $m<\mfix/\alpha$ and positive beyond.
(i) For $q>1$ the numerator grows like $m$ and the denominator like $m^q$, so
$\Theta_p(m)\to0$ as $m\to\infty$; a continuous function that is positive on
$(\mfix/\alpha,\infty)$ and tends to $0$ at both ends of that interval attains
its supremum. (ii) At $q=1$,
$\Theta_1'(m)$ has the sign of
$\alpha(c_P m+\Paux) - c_P(\alpha m-\mfix) = \alpha\Paux + c_P\mfix>0$, and
$\Theta_1(m)\to\alpha\eb/c_P$. The mass closing duration $t$ is
$m = (\mfix + \Paux t/\eb)/(\alpha - c_P t/\eb)$, positive and finite exactly
when $t<t_{\lim}$; with $\Paux=0$ this is \cref{eq:constdlfeas}. (iii) For
$q<1$ the numerator outgrows the denominator. (iv) By the envelope theorem,
$\partial_m\Theta_p(\hat m)=0$ gives
$\dd\ln t_f^{\max}/\dd p = \partial_p\ln\Theta_p(\hat m)
= -\partial_p \ln(c_P \hat m^q + \Paux)$, and with
$\partial_p\ln c_P = \tfrac12\ln m_0$ and $\partial_p q = -\tfrac12$ this is
$-\vartheta(\tfrac12\ln m_0 - \tfrac12\ln\hat m)$.
\end{proof}

\begin{remark}
With $\Paux=0$ the maximiser is explicit: $\hat m = q\,\mfix/[\alpha(q-1)]$,
which tends to infinity as $q\downarrow1$, reconciling (i) with (ii). If the
anchor is the fold mass of the fixed-area closure at the nominal duration $t_f$,
then $\sqrt{m_0} = 2\alpha/(3\beta)$ and \cref{eq:tlim} gives exactly
$t_{\lim} = \tfrac32 t_f$, independently of $\Paux$.
\end{remark}

\Cref{fig:theta} draws $\Theta_p$ for five exponents. The maxima of (i), the
unattained ceiling of (ii) and the unbounded growth of (iii) are all visible,
and the maxima lie to the right of the anchor, which by (iv) is why the wall
rises with $p$ here.

\begin{figure}[tbp]
\centering
% Theta_p(m) / t_f against m / m0 at the engine defaults, anchor m0 = fold
% mass of the fixed-area closure (gen_b.theta). For p < 1 the curve has an
% interior maximum (marked); at p = 1 it rises to alpha e_b / c_P = (3/2) t_f
% without reaching it; for p > 1 it grows without bound.
\begin{tikzpicture}
\begin{axis}[paperplot, height=6.4cm, xmode=log,
  xlabel={gross mass relative to the anchor, $m/m_0$},
  ylabel={longest closed mission $\Theta_p(m)/t_f$},
  xmin=0.3, xmax=300, ymin=0, ymax=2.1, unbounded coords=jump,
  legend style={at={(0.5,1.03)}, anchor=south, font=\scriptsize},
  legend columns=4]
  \addplot[cA] table[x=x, y=p0] {results/fig/b_theta.dat};
  \addlegendentry{$p=0$}
  \addplot[cC] table[x=x, y=p1] {results/fig/b_theta.dat};
  \addlegendentry{$p=0.5$}
  \addplot[cE] table[x=x, y=p2] {results/fig/b_theta.dat};
  \addlegendentry{$p=0.75$}
  \addplot[cB] table[x=x, y=p3] {results/fig/b_theta.dat};
  \addlegendentry{$p=1$}
  \addplot[cD] table[x=x, y=p4] {results/fig/b_theta.dat};
  \addlegendentry{$p=1.25$}
  \addplot[cB, dashed, line width=0.7pt] table[x=x, y=t] {results/fig/b_theta_lim.dat};
  \addlegendentry{$\alpha\eb/(c_Pt_f)=\tfrac32$}
  \addplot[only marks, mark=*, mark size=2pt, black] table[x=x, y=t]
    {results/fig/b_theta_peaks.dat};
  \addlegendentry{$\hat m(p)$, $t_f^{\max}(p)$}
  \draw[construct] (axis cs:1,0) -- (axis cs:1,2.1);
  \node[lbl, cG, anchor=south west, rotate=90] at (axis cs:0.97,0.05) {anchor $m_0$};
  \draw[cG, densely dotted] (axis cs:0.3,1) -- (axis cs:300,1);
  \node[lbl, cG, anchor=south east] at (axis cs:290,1.0) {nominal $t_f$};
\end{axis}
\end{tikzpicture}

\caption{The longest mission a mass $m$ can close, $\Theta_p(m)$ of
\cref{eq:theta}, relative to the nominal $t_f$, at the engine's default
parameters with the anchor $m_0$ at the fold mass of the fixed-area closure.
For $p<1$ each curve has an interior maximum $(\hat m,t_f^{\max})$ (dots), and
these lie to the right of the anchor. At $p=1$ the curve rises towards
$\alpha\eb/c_P=\tfrac32t_f$ (dashed) without reaching it; at $p=1.25$ it grows
without bound. No fold and no ceiling are different statements: $p=1$ has the
first property and not the second.}
\label{fig:theta}
\end{figure}

\Cref{tab:pwall} records the computed values and \cref{fig:pwall} draws them
on a finer grid. The default parameters do not close at fixed area at the
nominal hour, so the common anchor is the fold mass, and the $p=1$ entry is
$\tfrac32 t_f$ as the remark predicts.

\begin{figure}[tbp]
\centering
% Endurance wall against the disk-area exponent (gen_b.pwall), same
% parameters and anchor as the table. (a) t_max / t_f: the curve is the
% direct maximisation of Theta_p, the crosses the bisection of
% sweep.endurance_wall; the open circle at p = 1 is the unattained
% supremum 3/2. (b) the wall mass relative to the anchor; the wall rises
% with p exactly where this exceeds one.
\begin{tikzpicture}
\begin{groupplot}[group style={group size=2 by 1, horizontal sep=1.6cm},
  halfplot, height=5.8cm, width=0.46\linewidth,
  xlabel={disk-area exponent $p$}, xmin=-0.25, xmax=1.5]
\nextgroupplot[ylabel={$t_f^{\max}/t_f$}, ymin=0.8, ymax=1.75,
  title={(a) the endurance wall}, unbounded coords=discard]
  \fill[cD!8] (axis cs:1,0.8) rectangle (axis cs:1.5,1.75);
  \node[lbl, cD, align=center] at (axis cs:1.25,1.25) {no wall:\\$\Theta_p\to\infty$};
  \addplot[cA] table[x=p, y=tdirect] {results/fig/b_pwall.dat};
  \addplot[only marks, mark=x, mark size=2.2pt, cB] table[x=p, y=t]
    {results/fig/b_pwall.dat};
  \addplot[only marks, mark=o, mark size=2.6pt, cA, line width=1pt] table[x=p, y=t]
    {results/fig/b_pwall_lim.dat};
  \draw[cB, dashed] (axis cs:0.55,1.5) -- (axis cs:1,1.5);
  \node[lbl, anchor=south east] at (axis cs:0.98,1.52) {$t_{\lim}=\tfrac32t_f$};
\nextgroupplot[ylabel={$\hat m/m_0$}, ymode=log, ymin=0.5, ymax=2000,
  title={(b) the mass at the wall}, xmax=1.0]
  \addplot[cA] table[x=p, y=x] {results/fig/b_pwall.dat};
  \draw[cD, dashed] (axis cs:-0.25,1) -- (axis cs:1,1);
  \node[lbl, cD, anchor=south east] at (axis cs:0.95,1.05) {$\hat m=m_0$};
  \node[lbl, anchor=north west, align=left] at (axis cs:-0.2,900)
    {wall rises with $p$\\where $\hat m>m_0$};
\end{groupplot}
\end{tikzpicture}

\caption{The endurance wall of \cref{prop:twall} on a fine grid of exponents,
at the parameters and anchor of \cref{tab:pwall}. (a) $t_f^{\max}/t_f$ from
direct maximisation of \cref{eq:theta} (line) and from the bisection of
\code{sweep.endurance\_wall} (crosses), which agree wherever the bisection
runs; within about $0.01$ of $p=1$ the bisection overflows in the fold mass
$(\alpha/q\beta)^{1/(q-1)}$ and only the direct value is shown. The circle at
$p=1$ is the unattained supremum $\tfrac32t_f$ of \cref{eq:tlim}; for $p>1$
there is no wall. (b) The mass at the wall exceeds the anchor for every $p<1$
and diverges as $p\uparrow1$, so by \cref{eq:wallslope} the wall rises
monotonically on this family.}
\label{fig:pwall}
\end{figure}

\begin{table}[t]
\centering
\caption{Endurance wall against the disk-area exponent at the engine's default
parameters ($t_f=\SI{\rPwallTf}{\second}$), anchored at the fold mass
$m_0=\SI{\rPwallAnchor}{\kilo\gram}$ of the fixed-area closure. Computed by bisection on
$t_f$ with \code{sweep.endurance\_wall} and checked against direct maximisation
of \cref{eq:theta}. The last column is $\hat m/m_0$; it exceeds one throughout,
so by \cref{eq:wallslope} the wall rises with $p$ for this anchor. At $p=1$ there
is no fold but there is still a ceiling, the unattained supremum
\cref{eq:tlim}.}
\label{tab:pwall}
\begin{tabular}{@{}rrlrr@{}}
\toprule
$p$ & $q=(3-p)/2$ & fold exists & $t_f^{\max}$ & $\hat m/m_0$ \\
\midrule
\inputrows{results/tab_pwall_rows}
\end{tabular}
\end{table}
